% GENERATED FILE. Edit canonical lessons, then run npm run generate:books.
% canonical-source-hash: fnv1a64:75e51315ff80d1b4
% canonical-lessons: BN-C529-monopputo, BN-C529-dotana, BN-C529-ruchi, BN-C529-jhoka, BN-C529-morji

\chapter{Torn Between Two}
\label{ch:bn-torn-between-two}
\hlchaptermodality{529}

\begin{chapteropening}
\textbf{By the end of this chapter:} \emph{I can say that I am torn, what I like and which way I lean.}

Complete the last lesson of chapter 529: I can say that I am torn, what I like and which way I lean.
\end{chapteropening}

\section[\texorpdfstring{\bn{মনঃপূত} (mônôppūtô)}{মনঃপূত (mônôppūtô)}]{\bn{মনঃপূত} (mônôppūtô) — pleasing; to one's liking}

\label{lesson:BN-C529-monopputo}



\begin{quote}
Before the new one: say the Bengali for to suffice, then the Bengali for to scrape together.
\end{quote}

\subsection*{You'll want to know: \bn{মনঃপূত}}
\textbf{\bn{মনঃপূত}} — \emph{mônôppūtô} — \textquotedblleft{}pleasing; to one's liking\textquotedblright{}.

\begin{quote}
\bn{জামাটা} \bn{মনঃপূত} \bn{হলে} \bn{আমি} \bn{সেটা} \bn{কিনব।}
\end{quote}

\begin{quote}
\emph{Jāmāṭā mônôppūtô hôle āmi seṭā kinbô.}
\end{quote}

\begin{quote}
If the shirt is to my liking, I'll buy it.
\end{quote}

\begin{cousinweb}
\textbf{\bn{মনঃপূত}} (\emph{mônôppūtô}) is Sanskrit \emph{manaḥpūta}, \textquotedblleft{}pure in mind\textquotedblright{} (Monier-Williams): \emph{manas}, \textquotedblleft{}mind\textquotedblright{}, and \emph{pūta}, \textquotedblleft{}purified\textquotedblright{}; the sense \textquotedblleft{}to one's liking\textquotedblright{} is Bengali's own.
\end{cousinweb}

\begin{grammarlens}[title={What you've built}]
One more word for saying what happens if.
\end{grammarlens}

\subsection*{Your turn}
\begin{itemize}
\raggedright
  \item \emph{Say it:} \emph{mônôppūtô}
  \item \emph{Say it:} \emph{mônôppūtô}, once more
  \item \emph{Say it:} \emph{Jāmāṭā mônôppūtô hôle āmi seṭā kinbô.} — If the shirt is to my liking, I'll buy it.
  \item \emph{Recall:} say the Bengali for to suffice, then the Bengali for to scrape together, then say \emph{mônôppūtô} again
\end{itemize}

\begin{tcolorbox}[breakable,colback=teal!4,colframe=teal!35!black,title={Before you move on}]
What does \bn{মনঃপূত} mean? (\textquotedblleft{}Pleasing; to one's liking\textquotedblright{}.) Say it once more.
\end{tcolorbox}
\section[\texorpdfstring{\bn{দোটানা} (doṭānā)}{দোটানা (doṭānā)}]{\bn{দোটানা} (doṭānā) — a dilemma; being torn two ways}

\label{lesson:BN-C529-dotana}



\begin{quote}
Before the new one: say the Bengali for to scrape together, then the Bengali for pleasing to oneself.
\end{quote}

\subsection*{You'll want to know: \bn{দোটানা}}
\textbf{\bn{দোটানা}} — \emph{doṭānā} — \textquotedblleft{}a dilemma; being torn two ways\textquotedblright{}.

\begin{quote}
\bn{আমি} \bn{দোটানায়} \bn{আছি}: \bn{শহরে} \bn{থাকব}, \bn{না} \bn{গ্রামে} \bn{ফিরব}?
\end{quote}

\begin{quote}
\emph{Āmi doṭānāy āchhi: shôhôre thākbô, nā grāme phirbô?}
\end{quote}

\begin{quote}
I'm torn: shall I stay in the city or go back to the village?
\end{quote}

\begin{cousinweb}
\textbf{\bn{দোটানা}} (\emph{doṭānā}) joins \bn{দো}, \textquotedblleft{}two\textquotedblright{}, usually traced to Persian \emph{do}, to \bn{টানা}, \textquotedblleft{}a pull\textquotedblright{}: being pulled two ways at once.
\end{cousinweb}

\begin{grammarlens}[title={What you've built}]
One more word for saying what happens if.
\end{grammarlens}

\subsection*{Your turn}
\begin{itemize}
\raggedright
  \item \emph{Say it:} \emph{doṭānā}
  \item \emph{Say it:} \emph{doṭānā}, once more
  \item \emph{Say it:} \emph{Āmi doṭānāy āchhi: shôhôre thākbô, nā grāme phirbô?} — I'm torn: shall I stay in the city or go back to the village?
  \item \emph{Recall:} say the Bengali for to scrape together, then the Bengali for pleasing to oneself, then say \emph{doṭānā} again
\end{itemize}

\begin{tcolorbox}[breakable,colback=teal!4,colframe=teal!35!black,title={Before you move on}]
What does \bn{দোটানা} mean? (\textquotedblleft{}A dilemma; being torn two ways\textquotedblright{}.) Say it once more.
\end{tcolorbox}
\section[\texorpdfstring{\bn{রুচি} (ruchi)}{রুচি (ruchi)}]{\bn{রুচি} (ruchi) — personal taste}

\label{lesson:BN-C529-ruchi}



\begin{quote}
Before the new one: say the Bengali for pleasing to oneself, then the Bengali for a dilemma.
\end{quote}

\subsection*{You'll want to know: \bn{রুচি}}
\textbf{\bn{রুচি}} — \emph{ruchi} — \textquotedblleft{}personal taste\textquotedblright{}.

\begin{quote}
\bn{আমার} \bn{আর} \bn{আমার} \bn{বোনের} \bn{রুচি} \bn{আলাদা।}
\end{quote}

\begin{quote}
\emph{Āmār ār āmār boner ruchi ālādā.}
\end{quote}

\begin{quote}
My sister and I have different tastes.
\end{quote}

\begin{cousinweb}
\textbf{\bn{রুচি}} (\emph{ruchi}) is a tatsama, Sanskrit \emph{ruci}, whose senses in Monier-Williams run from \textquotedblleft{}light, lustre\textquotedblright{} to \textquotedblleft{}liking, taste\textquotedblright{}.
\end{cousinweb}

\begin{grammarlens}[title={What you've built}]
One more word for saying what happens if.
\end{grammarlens}

\subsection*{Your turn}
\begin{itemize}
\raggedright
  \item \emph{Say it:} \emph{ruchi}
  \item \emph{Say it:} \emph{ruchi}, once more
  \item \emph{Say it:} \emph{Āmār ār āmār boner ruchi ālādā.} — My sister and I have different tastes.
  \item \emph{Recall:} say the Bengali for pleasing to oneself, then the Bengali for a dilemma, then say \emph{ruchi} again
\end{itemize}

\begin{tcolorbox}[breakable,colback=teal!4,colframe=teal!35!black,title={Before you move on}]
What does \bn{রুচি} mean? (\textquotedblleft{}Personal taste\textquotedblright{}.) Say it once more.
\end{tcolorbox}
\section[\texorpdfstring{\bn{ঝোঁকা} (jhõkā)}{ঝোঁকা (jhõkā)}]{\bn{ঝোঁকা} (jhõkā) — to lean; to stoop}

\label{lesson:BN-C529-jhoka}



\begin{quote}
Before the new one: say the Bengali for a dilemma, then the Bengali for personal taste.
\end{quote}

\subsection*{You'll want to know: \bn{ঝোঁকা}}
\textbf{\bn{ঝোঁকা}} — \emph{jhõkā} — \textquotedblleft{}to lean; to stoop\textquotedblright{}.

\begin{quote}
\bn{সামনে} \bn{বেশি} \bn{ঝুঁকলে} \bn{তুমি} \bn{পড়ে} \bn{যাবে।}
\end{quote}

\begin{quote}
\emph{Sāmne beshi jhũkle tumi pôṛe jābe.}
\end{quote}

\begin{quote}
If you lean too far forward, you'll fall.
\end{quote}

Past (he/she …): \bn{তিনি} \bn{সামনে} \bn{ঝুঁকলেন}, \emph{tini sāmne jhũklen}.

The stem vowel rises in the past: \bn{ঝুঁকলেন}, he leaned.

\begin{cousinweb}
\textbf{\bn{ঝোঁকা}} (\emph{jhõkā}) comes from a reconstructed \emph{jhukkati}, \textquotedblleft{}stoops\textquotedblright{}; CDIAL lists Bengali \emph{jhũkā}, \textquotedblleft{}to stoop\textquotedblright{} (CDIAL 5399). \bn{ঝুঁকি}, a risk, is probably from the same root.
\end{cousinweb}

\begin{grammarlens}[title={What you've built}]
One more word for saying what happens if.
\end{grammarlens}

\subsection*{Your turn}
\begin{itemize}
\raggedright
  \item \emph{Say it:} \emph{jhõkā}
  \item \emph{Say it:} \emph{jhõkā}, once more
  \item \emph{Say it:} \emph{Sāmne beshi jhũkle tumi pôṛe jābe.} — If you lean too far forward, you'll fall.
  \item \emph{Recall:} say the Bengali for a dilemma, then the Bengali for personal taste, then say \emph{jhõkā} again
\end{itemize}

\begin{tcolorbox}[breakable,colback=teal!4,colframe=teal!35!black,title={Before you move on}]
What does \bn{ঝোঁকা} mean? (\textquotedblleft{}To lean; to stoop\textquotedblright{}.) Say it once more.
\end{tcolorbox}
\section[\texorpdfstring{\bn{মর্জি} (môrji)}{মর্জি (môrji)}]{\bn{মর্জি} (môrji) — a whim; a mood}

\label{lesson:BN-C529-morji}



\begin{quote}
Before the new one: say the Bengali for personal taste, then the Bengali for to lean forward.
\end{quote}

\subsection*{You'll want to know: \bn{মর্জি}}
\textbf{\bn{মর্জি}} — \emph{môrji} — \textquotedblleft{}a whim; a mood\textquotedblright{}.

\begin{quote}
\bn{আমার} \bn{মর্জি} \bn{হলে} \bn{আমি} \bn{সারাদিন} \bn{বই} \bn{পড়ি।}
\end{quote}

\begin{quote}
\emph{Āmār môrji hôle āmi sārādin bôi pôṛi.}
\end{quote}

\begin{quote}
When the mood takes me, I read books all day.
\end{quote}

\begin{cousinweb}
\textbf{\bn{মর্জি}} (\emph{môrji}) is usually traced to Arabic \emph{marḍī}, \textquotedblleft{}what is pleasing\textquotedblright{}, through Persian.
\end{cousinweb}

\begin{grammarlens}[title={What you've built}]
One more word for saying what happens if.
\end{grammarlens}

\subsection*{Your turn}
\begin{itemize}
\raggedright
  \item \emph{Say it:} \emph{môrji}
  \item \emph{Say it:} \emph{môrji}, once more
  \item \emph{Say it:} \emph{Āmār môrji hôle āmi sārādin bôi pôṛi.} — When the mood takes me, I read books all day.
  \item \emph{Recall:} say the Bengali for personal taste, then the Bengali for to lean forward, then say \emph{môrji} again
\end{itemize}

\begin{tcolorbox}[breakable,colback=teal!4,colframe=teal!35!black,title={Before you move on}]
What does \bn{মর্জি} mean? (\textquotedblleft{}A whim; a mood\textquotedblright{}.) Say it once more.
\end{tcolorbox}
